% ส่วนเพิ่มเติมบทที่ 1 ต่อจากหัวข้อ 1.6
\section*{1.7 ที่มาและความสำคัญของการพัฒนานโยบายการตัดสินใจเพิ่มเติม}
การพัฒนาระบบเชื่อมโยงอัตลักษณ์ในระยะแรกของโครงงานมุ่งให้ระบบสามารถนำข้อมูลโปรไฟล์จากหลายแพลตฟอร์มมาเปรียบเทียบและให้คะแนนความเป็นบุคคลเดียวกันได้ งานที่พัฒนาเพิ่มเติมในระยะต่อมามุ่งตอบคำถามว่า เมื่อมีคะแนนแล้ว ระบบควรตัดสินใจอย่างไรจึงจะเหมาะสมกับการนำข้อมูลไปใช้ใน CRM เพราะการมีแบบจำลองที่แยกคู่เหมือนและคู่ต่างได้ดี ยังไม่เพียงพอที่จะกำหนดว่าคู่ใดควรถูกรวมทันที และคู่ใดควรรอการตรวจสอบ

หัวข้อ 1.7--1.13 จึงขยายรายละเอียดจากที่มา วัตถุประสงค์ ขอบเขต และวิธีดำเนินงานเดิม โดยเน้นส่วนที่โครงงานได้พัฒนาเพิ่ม คือการค้นหานโยบายการตัดสินใจด้วยขั้นตอนวิธีเชิงพันธุกรรม การใช้หลักฐานหลายด้านอย่างเหมาะสมเมื่อข้อมูลไม่ครบ และการประเมินผลที่คำนึงถึงทั้งคุณภาพการเชื่อมโยงและภาระงานตรวจสอบ

\subsection*{1.7.1 จากปัญหาข้อมูลซ้ำสู่ความเสี่ยงของการรวมข้อมูลผิด}
ข้อมูลลูกค้าที่กระจายอยู่หลายระบบทำให้การมองเห็นภาพรวมของลูกค้าไม่ครบถ้วน แต่การพยายามรวมข้อมูลให้ได้มากที่สุดก็อาจสร้างปัญหาใหม่ หากระบบนำโปรไฟล์ของคนละคนมาอยู่ภายใต้ตัวตนเดียวกัน ข้อมูลความสนใจ ช่องทางติดต่อ หรือประวัติที่นำไปใช้ประกอบการทำงานอาจปะปนกัน ผลของการรวมผิดจึงไม่ได้จำกัดอยู่ที่คำตอบหนึ่งแถวในตาราง แต่สามารถส่งต่อไปยังกระบวนการที่ใช้โปรไฟล์รวมนั้นด้วย

ในอีกด้านหนึ่ง หากระบบระมัดระวังมากจนปฏิเสธคู่ที่เป็นคนเดียวกัน ข้อมูลของลูกค้าก็ยังคงแยกอยู่หลายรายการ การวิเคราะห์หรือจัดลำดับลูกค้าอาจใช้ข้อมูลเพียงบางส่วน การเชื่อมโยงจึงต้องพิจารณาสมดุลระหว่างการป้องกันการรวมผิดกับการลดคู่ที่ตกหล่น โดยไม่มีเหตุผลที่จะสมมติว่าความผิดพลาดสองประเภทสร้างภาระเท่ากันในทุกบริบท

ตัวอย่างเช่น คู่โปรไฟล์สองคู่ได้รับคะแนนจากแบบจำลองใกล้เคียงกัน คู่แรกมีชื่อคล้ายกันและมีลิงก์อ้างอิงสอดคล้องกัน ส่วนคู่ที่สองมีเพียงชื่อที่คล้ายกันและไม่มีหลักฐานอื่น การใช้จุดตัดเดียวกับคะแนนเพียงค่าเดียวอาจให้คำตัดสินเหมือนกันทั้งสองคู่ แม้ข้อมูลประกอบต่างกัน ตัวอย่างนี้เป็นเหตุผลในการศึกษาว่าการนำหลักฐานเพิ่มเติมเข้ามาในชั้นการตัดสินใจช่วยให้ระบบจัดการกรณีก้ำกึ่งได้เหมาะสมขึ้นหรือไม่

\subsection*{1.7.2 ข้อจำกัดของการกำหนดเกณฑ์ด้วยค่าคงที่}
ระบบเดิมใช้คะแนนและเกณฑ์สองระดับเพื่อแยกคู่เป็น MATCH, REVIEW และ NO\_MATCH แนวทางนี้ช่วยให้ไม่ต้องบังคับตัดสินทุกคู่เป็นใช่หรือไม่ใช่ทันที แต่ค่าจุดตัดยังมีผลโดยตรงต่อจำนวนคู่ที่ระบบรวมเองและจำนวนคู่ที่ต้องส่งให้คนตรวจสอบ หากตั้งเกณฑ์ MATCH สูงขึ้น จำนวนการรวมอัตโนมัติอาจลดลง ขณะที่คู่ส่วนหนึ่งย้ายไปอยู่ในคิว REVIEW หากลดเกณฑ์ลง ระบบอาจจับคู่ได้มากขึ้นพร้อมกับความเสี่ยงที่จะรวมผิดเพิ่มขึ้น

ภายใต้นโยบายที่ใช้คะแนนเดียวและเกณฑ์สองระดับ ความกว้างของช่วง REVIEW กำหนดพื้นที่ที่ระบบยอมเว้นการตัดสินใจอัตโนมัติ การขยายช่วงนี้อาจเพิ่มโอกาสให้เจ้าหน้าที่ช่วยแก้กรณียาก แต่ต้องแลกกับเวลาทำงาน การลดจำนวน REVIEW จึงเป็นเป้าหมายที่ต้องอ่านร่วมกับผลของคู่ที่ย้ายออกจากคิว ไม่ใช่นับเพียงว่าคิวสั้นลงแล้วถือว่าระบบดีขึ้น

ด้วยเหตุนี้ ผู้วิจัยจึงเพิ่มขั้นตอนค้นหาเกณฑ์จากข้อมูล validation และกำหนดต้นทุนสำหรับความผิดพลาดกับการตรวจสอบไว้ล่วงหน้า การเปรียบเทียบนโยบายบนข้อมูลเดียวกันทำให้สามารถอธิบายได้ว่าผลที่เปลี่ยนแปลงเกิดจากการปรับวิธีตัดสินใจมากน้อยเพียงใด

\subsection*{1.7.3 เหตุผลที่ใช้ GA และหลักฐานหลายด้าน}
นโยบายที่พัฒนาระยะแรกเพิ่มการใช้ความคล้ายของชื่อช่วยตัดสินคู่ในช่วง REVIEW ทำให้การค้นหามีพารามิเตอร์มากกว่าจุดตัดของคะแนนโมเดลเพียงสองค่า ต่อมาเมื่อขยายไปใช้หลักฐานหลายด้าน จำนวนพารามิเตอร์เพิ่มเป็นเกณฑ์สองระดับและน้ำหนักแปดค่า การเลือกค่าทีละตัวด้วยมืออาจไม่สะท้อนผลที่แต่ละค่ามีร่วมกัน เช่น การเพิ่มน้ำหนักชื่ออาจต้องเปลี่ยนเกณฑ์ MATCH จึงจะรักษาระดับความแม่นยำเดิมได้

GA เป็นทางเลือกหนึ่งสำหรับการค้นหาชุดพารามิเตอร์ร่วมกัน โดยสร้างนโยบายหลายชุด ประเมินผล แล้วนำชุดที่เหมาะสมไปสร้างคำตอบในรุ่นถัดไป การใช้ GA กับงานเชื่อมโยงข้อมูลมีตัวอย่างในงานของ Sun และคณะ [11] ซึ่งศึกษาการเรียนรู้กฎจับคู่ร่วมกับ active learning สำหรับโครงงานนี้ ผู้วิจัยประยุกต์ GA กับชั้นการตัดสินใจหลังได้คะแนนจากแบบจำลอง โดยตรึงคะแนนไว้ระหว่างเปรียบเทียบนโยบาย

อีกประเด็นที่ต้องพิจารณาคือความพร้อมของข้อมูล โปรไฟล์บางคู่มีประวัติย่อ ลิงก์ หรือภาพ ขณะที่บางคู่ไม่มี การไม่มีข้อมูลไม่ได้ให้ความหมายเดียวกับการมีข้อมูลที่ขัดแย้งกัน ผู้วิจัยจึงเพิ่มตัวบ่งชี้ความพร้อมของหลักฐาน และคำนวณคะแนนถ่วงน้ำหนักจากช่องที่มีข้อมูลจริงในแต่ละคู่ เพื่อศึกษานโยบายที่รองรับสภาพข้อมูลไม่ครบได้อย่างชัดเจน

\section*{1.8 วัตถุประสงค์ของงานที่พัฒนาเพิ่มเติม}
วัตถุประสงค์ในส่วนนี้ขยายจากเป้าหมายการเชื่อมโยงข้อมูลและเตรียมข้อมูลสำหรับ CRM ที่กล่าวไว้ในหัวข้อ 1.2 โดยมุ่งให้แต่ละวัตถุประสงค์มีผลลัพธ์ที่ตรวจสอบได้จากการทดลอง ดังนี้

\subsection*{1.8.1 พัฒนานโยบายการตัดสินใจที่ค้นหาพารามิเตอร์จากข้อมูล}
เพื่อพัฒนาวิธีเลือกเกณฑ์สำหรับ MATCH, REVIEW และ NO\_MATCH ด้วย GA โดยใช้ข้อมูล validation เป็นแหล่งเลือกพารามิเตอร์ และตรึงนโยบายก่อนประเมินกับ test ผลงานที่ต้องได้คือชุดพารามิเตอร์ที่นำกลับมาใช้ได้ พร้อมข้อมูลการตั้งค่าและผล validation ที่ใช้ประกอบการเลือก

\subsection*{1.8.2 ขยายการใช้หลักฐานในชั้นการตัดสินใจ}
เพื่อพัฒนานโยบายที่รวมคะแนนโมเดลกับหลักฐานด้านชื่อ ประวัติย่อ ลิงก์ ตำแหน่งที่อยู่ สัญญาณทางสังคม ภาพ และรูปแบบการเขียนตามข้อมูลที่มี โดยแยกค่าความคล้ายออกจากสถานะการมีข้อมูล ผลงานที่ต้องได้คือตารางหลักฐานรายคู่และสูตรการคำนวณที่ไม่ลดคะแนนเพียงเพราะช่องข้อมูลขาดหาย

\subsection*{1.8.3 เปรียบเทียบผลด้านคุณภาพและภาระงานอย่างเป็นระบบ}
เพื่อเปรียบเทียบระบบเดิม นโยบายคะแนนอย่างเดียว นโยบายใช้ชื่อช่วยตัดสิน และนโยบายแบบถ่วงน้ำหนัก ภายใต้คู่ข้อมูลและนิยามต้นทุนเดียวกัน โดยพิจารณา precision, recall, F1 จำนวนการรวมผิด จำนวนคู่ที่ปฏิเสธผิด และจำนวน REVIEW ร่วมกัน

\subsection*{1.8.4 จัดทำผลลัพธ์ที่ตรวจสอบย้อนกลับและส่งต่อได้}
เพื่อบันทึกแหล่งข้อมูล รุ่นนโยบาย ค่าเกณฑ์ น้ำหนัก และคำตัดสินรายคู่ให้สัมพันธ์กัน รองรับการตรวจสอบความแตกต่างจากระบบเดิม และเตรียมผลสำหรับกระบวนการรวมโปรไฟล์หรือจัดคิวตรวจสอบใน CRM โดยไม่ถือว่าการส่งออกผลเท่ากับการนำไปใช้งานจริงครบวงจรแล้ว

\section*{1.9 คำถามการศึกษาและแนวทางหาคำตอบ}
คำถามของงานเพิ่มเติมมุ่งตรวจสอบผลของการออกแบบนโยบายภายใต้ข้อมูลและข้อจำกัดที่โครงงานมีอยู่ การศึกษาไม่ได้ตั้งข้อสรุปล่วงหน้าว่า GA ต้องดีกว่าทุกแนวทาง หรือการเพิ่มหลักฐานต้องทำให้ทุกตัวชี้วัดสูงขึ้น

\begin{center}\small
\begin{tabularx}{\linewidth}{p{1.1cm}XX}
\toprule
คำถาม & ประเด็นที่ต้องการศึกษา & วิธีตรวจสอบ \\
\midrule
RQ1 & การค้นหาเกณฑ์ด้วย GA เปลี่ยนต้นทุนและภาระ REVIEW จากระบบเดิมอย่างไร & เปรียบเทียบระบบเดิมกับนโยบายคะแนนอย่างเดียวบนคู่ชุดเดียวกัน \\
RQ2 & การรวมหลักฐานหลายด้านให้ผลต่างจากการใช้ชื่อช่วยตัดสินอย่างไร & เปรียบเทียบนโยบายสี่ค่ากับนโยบายสิบค่า โดยตรึงคะแนนโมเดล \\
RQ3 & ความพร้อมของข้อมูลส่งผลต่อการใช้น้ำหนักและการอ่านผลอย่างไร & รายงาน availability, coverage และน้ำหนักของกลุ่มที่เปิดใช้จริง \\
RQ4 & นโยบายที่เลือกจาก validation ให้ผลบน test และข้อมูลรวมต่างกันอย่างไร & ตรึงพารามิเตอร์ก่อนทดสอบ และรายงานแต่ละกลุ่มข้อมูลแยกกัน \\
\bottomrule
\end{tabularx}
\par ตารางที่ 1.2 คำถามการศึกษาของงานที่พัฒนาเพิ่มเติม
\end{center}

ผลที่ใช้ตอบคำถามควรแสดงการแลกเปลี่ยนระหว่างตัวชี้วัดอย่างตรงไปตรงมา เช่น หากต้นทุนลดลงแต่ recall ลดลงด้วย ต้องอธิบายว่าสัดส่วนคู่ที่ย้ายไปเป็น NO\_MATCH หรือ REVIEW เปลี่ยนอย่างไร การเปลี่ยนแปลงดังกล่าวเป็นข้อมูลสำหรับเลือกนโยบายตามบริบท มากกว่าจะใช้ประกาศว่านโยบายหนึ่งเหมาะสมกับทุกสถานการณ์

\section*{1.10 ขอบเขตและเงื่อนไขของงานเพิ่มเติม}
\subsection*{1.10.1 ขอบเขตข้อมูลและหน่วยที่ใช้ศึกษา}
งานใช้ข้อมูลโปรไฟล์ข้ามแพลตฟอร์มที่ผ่านกระบวนการจัดเตรียมในโครงงาน โดยมีแหล่งตั้งต้นสัมพันธ์กับ Dataset-LinkSocial [7] ซึ่งเผยแพร่ข้อมูลจาก Twitter, Instagram และ Google+ หน่วยที่นโยบายตัดสินคือคู่โปรไฟล์ ไม่ใช่ลูกค้าหนึ่งรายหรือยอดขายหนึ่งรายการ จำนวนแถวในข้อมูลที่ประมวลผลแล้วจึงต้องอ้างอิงไฟล์ของรอบทดลอง และแยกจากจำนวนที่รายงานในคลังข้อมูลสาธารณะ

ในรอบข้อมูลเต็มที่ใช้เป็นหลักฐานของงานเพิ่มเติม มีคู่ 2,086,245 คู่ ประกอบด้วยคู่ที่ใช้คะแนนโมเดล 2,073,842 คู่ และคู่จาก Exact-First 12,403 คู่ การค้นหานโยบายใช้เฉพาะคู่ validation ในส่วนที่ผ่านการให้คะแนนโมเดล ส่วนคู่ Exact-First คงคำตัดสินเดิมไว้เพื่อให้ขอบเขตการเปลี่ยนแปลงชัดเจน

\subsection*{1.10.2 ขอบเขตของแบบจำลองและนโยบาย}
การทดลองนโยบายเปรียบเทียบวิธีนำคะแนนไปตัดสินใจ โดยไม่ได้ฝึกแบบจำลองใหม่ระหว่างการเปรียบเทียบแต่ละนโยบายในรอบเดียวกัน โครงสร้างแบบถ่วงน้ำหนักรองรับข้อมูลแปดช่อง แต่ช่องที่ไม่มีหลักฐานพร้อมใช้จะไม่ถูกนำมาค้นหาน้ำหนัก ในรอบข้อมูลเต็มที่รายงาน location และ social ไม่ได้ใช้งานตามสภาพของตารางข้อมูล จึงไม่อ้างว่าเป็นการใช้ครบทุกด้าน

การศึกษานี้ไม่ได้เพิ่มคู่ผู้สมัครที่หลุดจากขั้น candidate generation กลับเข้ามา ผลของ GA จึงอยู่ภายในขอบเขตคู่ที่ระบบส่งมาให้ตัดสิน ความครอบคลุมของขั้นสร้างคู่ยังเป็นข้อจำกัดที่ต้องรายงานแยกจากความสามารถของนโยบาย

\subsection*{1.10.3 ขอบเขตการประเมินต้นทุน}
ต้นทุนกำหนดจากสมมติฐานเวลาทำงานและค่าแรงของโครงงาน โดยให้การรวมผิดมีน้ำหนัก 2 การปฏิเสธคู่บวกมีน้ำหนัก 1 และ REVIEW มีน้ำหนัก 0.1 ไม่รวมรายได้ที่สูญเสีย ผลกระทบต่อชื่อเสียง หรือค่าใช้จ่ายของการแก้ไขระบบที่ไม่ได้มีข้อมูลวัดจริง การเปรียบเทียบจึงใช้ตอบว่านโยบายใดมีต้นทุนต่ำกว่าภายใต้สมมติฐานเดียวกัน

การเลือกนโยบายใช้ validation ส่วน test ใช้ประเมินหลังตรึงค่า การรายงานผลบนคู่ทั้งหมดมีไว้แสดงปริมาณงานของข้อมูลชุดนั้น และไม่ใช้แทนผลบน test นอกจากนี้ การทดสอบหลาย seed เป็นการตรวจผลจากการสุ่มค้นหา ไม่เท่ากับการทดสอบกับข้อมูลจากหลายองค์กรหรือหลายช่วงเวลา

\subsection*{1.10.4 ขอบเขตการเชื่อมต่อกับงาน CRM}
ผลที่จัดเตรียมสำหรับกระบวนการปลายทาง ได้แก่ รหัสคู่ คะแนน คำตัดสิน และข้อมูลที่ใช้ติดตามรุ่นนโยบาย ส่วนแนวคิดด้านการจัดลำดับลูกค้า การวิเคราะห์เส้นทางลูกค้า และการแนะนำกิจกรรมยังเป็นบริบทของโครงงานภาพรวม การปรับ decision policy ครั้งนี้ไม่ได้ใช้ข้อมูลการซื้อจริงมาทดสอบว่าการรวมโปรไฟล์ทำให้อัตราการซื้อเพิ่มขึ้น

หากนำผลไปใช้สร้างโปรไฟล์รวม ต้องตรวจความถูกต้องในระดับกลุ่มบุคคลเพิ่มเติม เพราะการจับคู่หลายคู่ต่อกันอาจทำให้ความผิดพลาดเพียงคู่เดียวเชื่อมกลุ่มข้อมูลที่ไม่ควรอยู่ร่วมกันได้ ผลดีในระดับคู่จึงยังไม่เพียงพอสำหรับยืนยันความถูกต้องของโปรไฟล์รวมทุกกลุ่ม

\section*{1.11 ภาพรวมวิธีดำเนินงานและการเชื่อมโยงกับงานเดิม}
งานเพิ่มเติมเริ่มจากทบทวนคะแนนและคำตัดสินของระบบเดิม เพื่อกำหนดฐานเปรียบเทียบ จากนั้นตรวจรหัสคู่และการแบ่งข้อมูล สร้างตารางหลักฐานพร้อมตัวบ่งชี้ความพร้อมใช้งาน แล้วกำหนดนโยบายและต้นทุนที่ใช้เปรียบเทียบให้ตรงกัน การค้นหาด้วย GA ดำเนินการบน validation ก่อนบันทึกพารามิเตอร์ที่เลือกและนำไปประเมิน test

หลังทดสอบ ผู้วิจัยนำพารามิเตอร์เดิมไปคำนวณคำตัดสินของคู่ทั้งหมด คงผล Exact-First ไว้ และส่งออกผลแยกตามนโยบาย การวิเคราะห์ปลายทางพิจารณาจำนวนคู่แต่ละสถานะ ความผิดพลาด และหลักฐานที่ใช้ประกอบคำตัดสิน รายละเอียดสูตรและการตั้งค่าการทดลองอธิบายไว้ในหัวข้อ 3.6--3.8

\begin{center}\small
\begin{tabularx}{\linewidth}{p{3cm}XX}
\toprule
องค์ประกอบ & งานเดิมที่ใช้ต่อ & งานที่เพิ่มในระยะนี้ \\
\midrule
ข้อมูลและคู่ผู้สมัคร & โปรไฟล์ที่จัดเตรียมแล้ว รหัสคู่ และผลสร้างคู่ & ตรวจความตรงกันของคู่กับตารางหลักฐานสำหรับนโยบาย \\
คะแนน & คะแนนจากแบบจำลองและคำตัดสินเดิม & ตรึงคะแนนเพื่อเปรียบเทียบวิธีตัดสินใจ \\
หลักฐาน & คุณลักษณะที่คำนวณไว้ในโครงงาน & จัดกลุ่มคะแนนและสถานะความพร้อมใช้งาน \\
การตัดสินใจ & MATCH, REVIEW และ NO\_MATCH & ค้นหาเกณฑ์และน้ำหนักด้วย GA \\
การประเมิน & ตัวชี้วัดคุณภาพการเชื่อมโยง & เพิ่มต้นทุน ภาระตรวจสอบ และผลแยกขอบเขตข้อมูล \\
การส่งออก & ผลจับคู่สำหรับงานปลายทาง & บันทึกรุ่นนโยบายและผลเปรียบเทียบรายคู่ \\
\bottomrule
\end{tabularx}
\par ตารางที่ 1.3 ความเชื่อมโยงระหว่างงานเดิมกับงานเพิ่มเติม
\end{center}

\section*{1.12 ประโยชน์ที่คาดว่าจะได้รับและผลงานที่ส่งมอบ}
ประโยชน์หลักคือการทำให้เหตุผลในการเลือกนโยบายชัดเจนขึ้น ผู้ใช้ผลการทดลองสามารถเห็นว่าการตั้งเกณฑ์และน้ำหนักทำให้การรวมผิด คู่ตกหล่น และงานตรวจสอบเปลี่ยนไปอย่างไร แทนการพิจารณาคะแนนจำแนกเพียงค่าเดียว นอกจากนี้ การแยกข้อมูลขาดออกจากค่าความคล้ายต่ำช่วยให้ตีความผลของหลักฐานแต่ละด้านได้รอบคอบขึ้น

ในด้านการพัฒนาระบบ การตรึงและบันทึกนโยบายช่วยให้เปรียบเทียบผลระหว่างรุ่นได้ เมื่อมีข้อมูลหรือสมมติฐานต้นทุนใหม่ สามารถนำกระบวนการเดิมกลับมาค้นหาและประเมินอีกครั้งได้โดยมีฐานอ้างอิง ผลงานที่ส่งมอบประกอบด้วยตารางหลักฐานรายคู่ ข้อมูลความครอบคลุม ผลค้นหาจากหลาย seed ไฟล์นโยบายที่ตรึงไว้ ตารางตัวชี้วัด และคำตัดสินรายคู่

ประโยชน์ด้านการลดเวลาตรวจสอบและการปรับคุณภาพข้อมูล CRM เป็นสิ่งที่คาดหวังจากการออกแบบ การยืนยันผลในองค์กรจริงยังต้องเก็บข้อมูลการตรวจสอบ การแก้ไขการรวมผิด และผลต่อกระบวนการทำงานปลายทางเพิ่มเติม จึงไม่ใช้ผลของการทดลองออฟไลน์แทนหลักฐานประโยชน์ทางธุรกิจที่เกิดขึ้นจริง

\section*{1.13 ลำดับการดำเนินงานและนิยามที่ใช้ในส่วนเพิ่มเติม}
การดำเนินงานแบ่งเป็นช่วงต่อเนื่อง ได้แก่ ตรวจฐานข้อมูลและผลเดิม เตรียมหลักฐาน ออกแบบนโยบาย ค้นหาพารามิเตอร์ ตรึงและทดสอบนโยบาย และจัดทำผลเปรียบเทียบ แต่ละช่วงต้องได้ข้อมูลที่ใช้เป็นจุดเริ่มต้นของช่วงถัดไป เช่น ต้องตรวจรหัสคู่และความพร้อมของหลักฐานก่อนค้นหาน้ำหนัก และต้องบันทึกนโยบายก่อนประเมิน test การเรียงงานเช่นนี้ทำให้ติดตามได้ว่าความคลาดเคลื่อนเกิดขึ้นในขั้นใด

ในส่วนเพิ่มเติมนี้ ``คะแนนโมเดล'' หมายถึงค่าที่แบบจำลองให้กับคู่โปรไฟล์ ส่วน ``คะแนนนโยบาย'' หมายถึงค่าที่คำนวณจากหลักฐานและน้ำหนักเพื่อใช้ตัดสินใจ ``นโยบายที่ตรึงไว้'' หมายถึงชุดเกณฑ์และน้ำหนักที่เลือกเสร็จก่อนขั้นประเมิน test และ ``ข้อมูลพร้อมใช้งาน'' หมายถึงมีหลักฐานที่สามารถคำนวณและใช้ตามนิยามของช่องนั้นได้ ไม่ได้หมายความว่าหลักฐานนั้นยืนยันความเป็นบุคคลเดียวกันเสมอไป

การใช้คำเหล่านี้อย่างแยกกันช่วยเชื่อมบทนำกับทฤษฎีในบทที่ 2 และวิธีดำเนินงานในบทที่ 3 โดยเฉพาะการแยกคุณภาพของแบบจำลองออกจากผลของนโยบาย และแยกผลการทดลองออกจากเป้าหมายการใช้งานในภาพรวมของโครงงาน
